% ==============================================================================
% reports/chapters/05_thiet_ke_thuc_nghiem.tex
% CHƯƠNG 5: THIẾT KẾ THỰC NGHIỆM
% ==============================================================================

\chapter{Thiết Kế Thực Nghiệm}
\label{chap:thiet_ke}

\section{Môi Trường Thực Thi Phần Cứng và Phần Mềm}

Toàn bộ các thực nghiệm trong nghiên cứu này được triển khai trên môi trường máy trạm độc lập với các thông số kỹ thuật được kiểm soát chặt chẽ:
\begin{itemize}
    \item \textbf{Phần cứng}: Vi xử lý đa nhân x86\_64, bộ nhớ RAM 32 GB, hỗ trợ tập lệnh SIMD AVX2.
    \item \textbf{Hệ điều hành}: Windows 11 / Linux POSIX tương thích.
    \item \textbf{Nền tảng phần mềm}: Python 3.12 cùng hệ thống thư viện khoa học dữ liệu gồm \texttt{numpy==1.26.4}, \texttt{pandas==2.2.2}, \texttt{scikit-learn==1.6.1}, \texttt{xgboost==3.0.5}, \texttt{shap==0.48.0} và \texttt{matplotlib==3.9.2}.
\end{itemize}

\section{Kiểm Soát Tính Ngẫu Nhiên}

Để đảm bảo khả năng tái lập tuyệt đối (Reproducibility), toàn bộ các tiến trình sinh số giả ngẫu nhiên từ quá trình phân chia dữ liệu, lấy mẫu bootstrap của Random Forest, khởi tạo trọng số của XGBoost và lấy mẫu nền của SHAP đều được khóa cứng với một hạt giống ngẫu nhiên duy nhất:
\begin{equation}
    \texttt{random\_state} = 42
\end{equation}

\section{Xử Lý Mất Cân Bằng Lớp Trên Tập Huấn Luyện}

Như đã phân tích tại Bảng \ref{tab:dataset_splits}, trong tập $\Dtrain$ theo thời gian, số lượng mẫu bình thường là \DataTrainBenign\ mẫu trong khi số lượng mẫu tấn công chỉ có \DataTrainAttack\ mẫu (tỷ lệ tấn công là \DataTrainAttackPercent). Tỷ lệ mất cân bằng lớp xấp xỉ $1 : 16.58$.

Để khắc phục hiện tượng mô hình bị thiên lệch về lớp đa số mà không làm biến dạng phân phối không gian mẫu (như các phương pháp tổng hợp dữ liệu SMOTE thường gây nhiễu ranh giới mạng), dự án áp dụng kỹ thuật gán trọng số lớp (Cost-sensitive Weighting) tính toán duy nhất từ $\Dtrain$:
\begin{itemize}
    \item \textbf{Decision Tree và Random Forest}:\\
    Sử dụng tham số trọng số \texttt{class\_weight='balanced'}:
    \begin{equation}
        w_c = \frac{N_{\text{train}}}{2 \cdot N_{\text{train}, c}} \implies w_0 = \frac{90,273}{2 \times 85,139} \approx 0.53, \quad w_1 = \frac{90,273}{2 \times 5,134} \approx 8.79
    \end{equation}
    \item \textbf{Đối với XGBoost}: Sử dụng tham số \texttt{scale\_pos\_weight}:
    \begin{equation}
        \text{scale\_pos\_weight} = \frac{N_{\text{train, normal}}}{N_{\text{train, attack}}} = \frac{85,139}{5,134} \approx 16.5834
    \end{equation}
\end{itemize}

\section{Thiết Lập Không Gian Tìm Kiếm Siêu Tham Số}

Để đảm bảo tính khách quan và khả năng tái lập tuyệt đối, không gian siêu tham số được chuẩn hóa thành lưới tìm kiếm chính thức (\textit{Canonical Benchmark Grid}) và được khóa cấu hình trong tệp lưu vết \texttt{artifacts/week3\_week4/frozen.json}. Các tham số quan trọng nhất của từng thuật toán được xác định như sau:

\begin{enumerate}
    \item \textbf{Cây Quyết Định (Decision Tree)}:
    \begin{itemize}
        \item Tiêu chuẩn phân tách (\texttt{criterion}): Mặc định \texttt{'gini'} (đã khảo sát bổ sung \texttt{'entropy'} trong các thử nghiệm thăm dò của TV2).
        \item Độ sâu tối đa (\texttt{max\_depth}): \texttt{[6, 8, 12, None]} (với \texttt{None} là cây phát triển hoàn toàn đến khi lá thuần nhất).
        \item Số mẫu tối thiểu tại nút lá (\texttt{min\_samples\_leaf}): \texttt{[5, 10]}.
    \end{itemize}
    \item \textbf{Rừng Ngẫu Nhiên (Random Forest)}:
    \begin{itemize}
        \item Số lượng cây (\texttt{n\_estimators}): \texttt{[100, 150, 300]}.
        \item Độ sâu tối đa (\texttt{max\_depth}): \texttt{[10, 16, None]}.
        \item Số mẫu tối thiểu tại nút lá (\texttt{min\_samples\_leaf}): \texttt{[5]}.
        \item Số đặc trưng tại mỗi nút (\texttt{max\_features}): \texttt{'sqrt'}.
    \end{itemize}
    \item \textbf{XGBoost}:
    \begin{itemize}
        \item Số lượng cây (\texttt{n\_estimators}): \texttt{[100, 200]}.
        \item Độ sâu cây (\texttt{max\_depth}): \texttt{[3, 4, 6]}.
        \item Tốc độ học (\texttt{learning\_rate}): \texttt{0.08}.
        \item Tỷ lệ lấy mẫu đặc trưng (\texttt{colsample\_bytree}): \texttt{0.8}.
        \item Tỷ lệ lấy mẫu dòng (\texttt{subsample}): \texttt{0.8}.
    \end{itemize}
\end{enumerate}

Bên cạnh lưới chính thức trên, các thực nghiệm mở rộng sử dụng \texttt{RandomizedSearchCV} (như tinh chỉnh 50 lượt lặp cho Decision Tree hoặc khảo sát tốc độ học $[0.05, 0.08, 0.1]$ cho XGBoost) được ghi nhận độc lập tại các thư mục thực nghiệm riêng lẻ để phục vụ phân tích độ nhạy. Chi tiết các cấu hình tối ưu được lựa chọn cho 12 mô hình được tổng hợp tại Phụ lục \ref{app:hyperparameters}.

\section{Tiêu Chí Lựa Chọn Mô Hình Chiến Thắng (Validation Selection)}

Trên tập kiểm định $\Dval$, các cấu hình được đánh giá theo thứ tự ưu tiên nghiêm ngặt:
\begin{enumerate}
    \item \textbf{Ưu tiên 1}: Điểm \textbf{F1-Score} trên lớp tấn công (tại ngưỡng mặc định $\tau = 0.5$).
    \item \textbf{Ưu tiên 2}: Điểm \textbf{Average Precision (AP)} trên đường cong Precision-Recall nếu F1 tương đương.
    \item \textbf{Ưu tiên 3}: Tỷ lệ báo động giả \textbf{FPR} thấp hơn nếu cả F1 và AP đều xấp xỉ nhau.
\end{enumerate}

Mô hình đạt điểm số tối ưu trên $\Dval$ sẽ được lựa chọn làm đại diện duy nhất của kịch bản đó để chuyển sang bước đánh giá trên tập kiểm thử $\Dtest$.

\section{Phân Định Rõ Hai Chế Độ Huấn Luyện (Model Regimes)}

Nhằm đảm bảo tính minh bạch khoa học và tránh việc diễn giải sai lệch kết quả, nghiên cứu phân định rõ ràng \textbf{hai chế độ thực nghiệm} độc lập đã được thực hiện:

\subsection{Chế độ 1: Canonical Strict Zero-Shot Regime (Chuẩn Mực Chính)}
Đây là giao thức đánh giá trung tâm của bài tập lớn được triển khai trong \path{scripts/run_remaining.py} và lưu trữ tại \path{artifacts/week3_week4/test_comparison.csv}:
\begin{itemize}
    \item Mô hình được \textbf{fit duy nhất trên tập $\Dtrain$} (chỉ chứa \FTP).
    \item Tập $\Dval$ chỉ được đưa vào phương thức \texttt{predict\_proba} để lựa chọn cấu hình tham số tốt nhất.
    \item Sau khi chọn được tham số, mô hình chiến thắng được huấn luyện lại \textbf{chỉ trên $\Dtrain$}, hoàn toàn không đụng tới $\Dval$.
    \item Đánh giá trên tập $\Dtest$ (chỉ chứa \SSH) $\rightarrow$ Đánh giá đúng năng lực Zero-Shot Subtype Transfer thực sự của mô hình.
\end{itemize}

\subsection{Chế độ 2: Subtype-Aware / Refit Regime (Case Study Mở Rộng)}
Đây là chế độ thực nghiệm phát sinh từ script tuning \texttt{src/ids/tune\_xgboost.py}:
\begin{itemize}
    \item Hàm \texttt{RandomizedSearchCV} sử dụng phân tách tùy biến trên tập gộp $\mathbf{X}_{\text{combined}} = \text{concat}([\mathbf{X}_{\text{train}}, \mathbf{X}_{\text{val}}])$.
    \item Do cờ \texttt{refit} mặc định của scikit-learn là \texttt{True}, sau khi tìm kiếm tham số tối ưu, thư viện đã \textbf{tự động huấn luyện lại mô hình tốt nhất trên toàn bộ tập gộp $\mathbf{X}_{\text{combined}}$}.
    \item Do $\mathbf{X}_{\text{val}}$ có chứa 1,886 mẫu \SSH, mô hình XGBoost trong chế độ này \textbf{đã được tiếp xúc và học trước các mẫu tấn công SSH} trước khi bước vào tập Test.
    \item Chế độ này được sử dụng như một \textbf{Case Study chuyên sâu về MLOps} trong Chương 7 để minh chứng cho tầm quan trọng của việc kiểm soát ranh giới dữ liệu và cờ refit trong NIDS.
\end{itemize}

\section{Đảm Bảo Tính Toàn Vẹn và Khả Năng Tái Lập}
Nhằm đảm bảo tính toàn vẹn khoa học và ngăn chặn tuyệt đối lỗi hồi quy trong toàn bộ đường ống xử lý dữ liệu, mọi ràng buộc về mốc thời gian, kỹ thuật Purge/Embargo và tính đóng băng của tập Test được tự động hóa kiểm định thông qua bộ kiểm thử hồi quy \texttt{tests/test\_protocol.py} (đạt 100\% tỷ lệ kiểm thử thành công, quy trình thực thi và mã nguồn chi tiết xem tại Phụ lục \ref{app:reproducibility}).
